%======================================================================
% INTEGRACIONES
%======================================================================

\chapter{Integraciones}

Sandra UI Consola proporciona múltiples opciones de integración con sistemas externos y herramientas de terceros.

\section{API REST}

La API REST permite integración programática con sistemas externos:

\begin{itemize}
    \item Endpoints para operaciones CRUD
    \item Autenticación mediante tokens JWT
    \item Formato JSON para intercambio de datos
    \item Documentación OpenAPI/Swagger
\end{itemize}

\section{Webhooks}

Los webhooks permiten recibir notificaciones en tiempo real:

\begin{itemize}
    \item Configuración de eventos triggering
    \item URLs de callback personalizables
    \item Reintentos automáticos en caso de fallo
    \item Firmas de verificación de payload
\end{itemize}

\section{Orquestación de Procesos (Fnx Engine)}

Sandra utiliza un motor de orquestación propietario para manejar tareas de larga duración:

\subsection{Gestión de PIDs}

Un PID (Process Identifier) es un número único asignado a cada proceso en ejecución, permitiendo su seguimiento y control individual.

El manejo de PIDs permite implementar cancellation, progreso tracking, y recuperación ante fallos de procesos de larga duración.

Cada instrucción pesada (escaneos de red, migraciones) se ejecuta en segundo plano con un PID único:

\begin{itemize}
\item \textbf{Ejecución Asíncrona}: El usuario no espera a que termine el proceso.
\item \textbf{Monitoreo en Tiempo Real}: Seguimiento del progreso del proceso.
\item \textbf{Cancelación}: Posibilidad de detener procesos si es necesario.
\item \textbf{Resultados Almacenados}: Los resultados se guardan para consulta posterior.
\end{itemize}

El patrón de ejecutar procesos pesados en background mejora significativamente la experiencia de usuario, evitando bloqueos de interfaz.

\subsection{Escaneo de Infraestructura}

Integración nativa con Nmap para descubrimiento automático de activos:

\begin{itemize}
\item \textbf{Descubrimiento de Dispositivos}: Identificación automática de servidores, routers, switches, dispositivos biométricos.
\item \textbf{Detección de Puertos}: Identificación de puertos abiertos en cada dispositivo.
\item \textbf{Análisis de Servicios}: Determinación de servicios activos en cada puerto.
\item \textbf{Exportación de Resultados}: Generación de reportes para auditoría.
\end{itemize}

La integración con Nmap permite mantener un inventario actualizado de activos de red, fundamental para la seguridad perimetral.

\clearpage
